\documentclass[10pt]{letter}
\usepackage[margin=0.65in]{geometry}
\usepackage{hyperref}
\usepackage{url}
\setlength{\parskip}{3pt}
\setlength{\parindent}{0pt}
\setlength{\topskip}{4pt}
\addtolength{\topmargin}{-0.3in}
\addtolength{\textheight}{80pt}

\signature{Lizhuo Zhang \\ zhanglizhuo@hunau.edu.cn}
\address{Lizhuo Zhang \\ Hunan Agricultural University \\ Changsha 410128, China}

\begin{document}

\begin{letter}{Editor-in-Chief \\ Engineering Applications of Artificial Intelligence}

\opening{Dear Editor-in-Chief,}

We submit the enclosed manuscript, ``AnchorScore: A CLIP-Based Diagnostic of MLLM Annotation Difficulty,'' as a Research Paper for consideration in Engineering Applications of Artificial Intelligence. Lizhuo Zhang (zhanglizhuo@hunau.edu.cn) is the corresponding author; Yan Ma has approved this submission.

The application context is an engineering-relevant and increasingly common one: multimodal large language models (MLLMs) are used to annotate classroom behavior data at industrial scale, yet their per-class accuracy varies dramatically (12\%--98\% in our setting) and measuring it costs roughly 14 hours of inference per model. This paper evaluates \emph{AnchorScore}, the per-class zero-shot accuracy of a frozen CLIP model, as a low-cost a priori difficulty diagnostic that flags the classes an MLLM is least likely to annotate reliably---for about 3 minutes of CLIP inference per dataset, roughly 1/270th of the compute of one MLLM pass. It is a resource-allocation tool for annotation pipelines, not an accuracy estimator.

The paper rests on four empirical findings. First, AnchorScore correlates with class-level MLLM accuracy on SCB5 (13 classes across three classroom sub-datasets, six MLLMs; Spearman $\rho{=}0.769$, $p{=}0.002$) and replicates on Stanford40 (40 classes; $\rho{=}0.817$, $p{<}0.001$). Second, the signal is absent from the obvious cheap alternatives---a 7B MLLM pilot on the same labeled images carries no class-difficulty signal at all ($\rho{=}{-}0.03$), CLIP's own prediction confidence routes no better than random, and DINOv2, ResNet-50, and MLLM self-verbalized uncertainty all fail as class-level predictors (the signal is only suggestively present for SigLIP, which the manuscript reads as a difference in proxy quality across model families)---establishing that the probe is specific rather than generic. Third, three engineering workflows translate the diagnostic into practice: hybrid CLIP/MLLM annotation routing (up to $+$23~pp over CLIP-only at about 44\% MLLM cost savings, with realized pipeline runs delivering $+$21.7~pp under the benchmark's canonical MLLM evaluation and $+$17.5~pp under the full-frame deployment condition where no annotation-derived bounding box is used), prompt disambiguation guided by CLIP confusion patterns, and review-priority prediction for human verification (AUC 0.84--0.94). Fourth, the signal attenuates on visually distant medical and satellite domains (pooled $\rho{=}0.462$), so the paper frames AnchorScore as a low-cost ranking diagnostic with documented operating limits rather than a universal estimator. The manuscript reports these limits in full, including a Mantel test showing that CLIP and MLLM confusion structures do not align.

We believe this work fits Engineering Applications of Artificial Intelligence: it delivers an AI-driven diagnostic embedded in a real annotation workflow, validated on public datasets (SCB5, Stanford40, MedMNIST, EuroSAT) with a public code repository, and it directly supports a practical engineering decision---where to spend limited MLLM annotation budget. It is an applied diagnostic contribution rather than a pure algorithmic or benchmark result.

This manuscript is a companion to our paper ``When Agreement Fails: Visual Anchoring and Multimodal Pseudo-Label Reliability in Classroom Behavior Recognition'' (Ma \& Zhang, IEEE Access, 2026), which is accepted and published. The two papers are complementary and contain no overlapping tables or figures: the companion paper provides the SCB5 pseudo-label reliability baseline---the data pipeline, confidence scoring, and inter-model agreement measures---and is retrospective (measuring reliability ex post). The present manuscript adds an \emph{a priori} diagnostic: AnchorScore correlation with per-class MLLM accuracy and three downstream workflows, i.e., it predicts where pseudo-labels fail before annotation. The AnchorScore construct corresponds to the companion paper's AnchorProxy (per-class zero-shot CLIP accuracy), independently recomputed under the full-frame deployment protocol rather than the bbox-cropped protocol; the manuscript states this relationship explicitly in its method section and cross-references the companion paper's per-class table.

This manuscript is original, has not been published elsewhere, and is not under consideration by another journal. All authors have read and approved the final version. The authors declare no competing interests, and this research received no external funding. SCB5 is public on Hugging Face; code, prompts, and result files are open at \url{https://github.com/zhanglizhuo/VisualAnchor}.

Thank you for your consideration.

\closing{Sincerely,}

\end{letter}

\end{document}